\documentclass[11pt,a4paper]{article}

\newcommand{\doctitle}{Fundamentos Técnicos}
\newcommand{\docsubtitle}{WARDEN - Wireless Attack Reproduction and Detection Environment}

\usepackage{warden-style}

\begin{document}

% -- Portada -------------------------------------------------------------------
\wardencoverpage{0.1}{Abril 2026}{Borrador}

% -- Historial de Revisiones ---------------------------------------------------
\begin{wardenrevisions}
\revrow{0.1}{Abril 2026}{Santiago Valencia León}{Borrador inicial}
\revrow{0.1}{Abril 2026}{Daniel Yadir Perea Murillo}{Borrador inicial}
\revrow{0.1}{Abril 2026}{Sofia Valdez Londono}{Borrador inicial}
\end{wardenrevisions}

% -- Tabla de Contenidos -------------------------------------------------------
\tableofcontents
\newpage

% ==============================================================================
\section{Introducción}
% ==============================================================================

\subsection{Propósito de este documento}
Este documento existe para que cualquier integrante del equipo WARDEN, sin
importar su experiencia previa en redes o ciberseguridad, pueda entender los
conceptos fundamentales que sustentan el proyecto. Está escrito de forma
narrativa: no es un glosario rápido sino una explicación construida pieza
por pieza, con analogías, para que al terminar de leerlo las palabras
\emph{beacon}, \emph{deauth}, \emph{modo monitor}, \emph{2.4 GHz}, o
\emph{handshake} dejen de ser jerga y pasen a ser ideas claras.

\subsection{Cómo leer este documento}
El documento está pensado para leerse \textbf{en orden}. Cada sección asume
que se entendió la anterior. La primera mitad construye el modelo mental
de cómo funciona WiFi en general, y la segunda mitad explica los ataques y
las técnicas que el proyecto WARDEN reproduce y detecta.

Si en algún momento aparece un término que no está claro, casi siempre es
porque se explica en una sección anterior. Volver a leer no es debilidad,
es la forma correcta de absorber este tipo de material.

\subsection{Audiencia}
Este documento está dirigido a los integrantes del equipo de desarrollo de
WARDEN. También sirve como material de respaldo durante la presentación
del proyecto: cualquier pregunta de fundamentos que el profesor evaluador
pueda hacer debería poder responderse con la información contenida aquí.

% ==============================================================================
\section{¿Qué es WiFi?}
% ==============================================================================

\subsection{Una conversación por radio}
Cuando una laptop, un celular o cualquier dispositivo ``se conecta a WiFi'',
lo que está ocurriendo es una conversación por ondas de radio entre el
dispositivo y un aparato llamado \textbf{punto de acceso} (en inglés
\emph{Access Point}, abreviado \textbf{AP}). El AP es típicamente un router
WiFi en una casa o empresa.

La idea es la misma que con una radio AM/FM: hay un emisor, hay un receptor,
hay ondas viajando por el aire. La diferencia clave es que en WiFi todos
los dispositivos pueden ser \emph{a la vez} emisores y receptores, y los
mensajes que se envían no son audio sino datos digitales: páginas web,
mensajes de WhatsApp, video de YouTube, etc.

\subsection{El nombre técnico: 802.11}
La radio comercial sigue estándares definidos por organismos internacionales
(qué frecuencias se usan, cómo se modula la señal). WiFi también. El estándar
que define cómo funciona WiFi se llama \textbf{IEEE 802.11}. La sigla IEEE es
del organismo que lo creó (Institute of Electrical and Electronics
Engineers), y 802.11 es el número del estándar.

Cuando alguien dice ``frame 802.11'' se refiere a un mensaje WiFi siguiendo
este estándar. WARDEN trabaja exclusivamente con frames 802.11.

\subsection{¿Por qué se llama WiFi entonces?}
``WiFi'' es un nombre comercial que la industria adoptó por ser más fácil de
recordar que ``IEEE 802.11''. Para fines del proyecto, ambos términos se
usan de forma intercambiable.

% ==============================================================================
\section{Espectro electromagnético y radiofrecuencia}
% ==============================================================================

\subsection{El espectro: de qué está hecha la luz}
La luz visible que perciben nuestros ojos, las microondas que calientan la
comida, los rayos X de un hospital, las ondas de radio AM/FM, y las señales
WiFi son todas el mismo fenómeno físico: \textbf{ondas electromagnéticas}.
Lo único que cambia entre ellas es su \emph{frecuencia}, es decir, qué tan
rápido oscilan.

El conjunto de todas las frecuencias posibles se llama \textbf{espectro
electromagnético}. Va desde frecuencias muy bajas (ondas de radio largas,
miles de kilómetros) hasta frecuencias muy altas (rayos gamma).

\subsection{Hertz, kilohertz, megahertz, gigahertz}
La frecuencia se mide en \textbf{hertz} (Hz). Un hertz significa ``una
oscilación por segundo''. Las cifras grandes se acortan con prefijos:

\begin{datatable}{L{2.5cm} L{4cm} L{8.8cm}}{Unidad & Valor & Ejemplo}
\drow{1 Hz   & 1 oscilación por segundo                & Latido del corazón en reposo (aprox.)}
\drow{1 kHz  & 1\,000 oscilaciones por segundo         & Sonido agudo audible}
\drow{1 MHz  & 1\,000\,000 oscilaciones por segundo    & Radio AM (alrededor de 1 MHz)}
\drow{1 GHz  & 1\,000\,000\,000 oscilaciones por segundo & WiFi 2.4 GHz, microondas}
\end{datatable}

WiFi opera en la región de los \textbf{gigahertz}: miles de millones de
oscilaciones por segundo. Específicamente, en dos bandas: alrededor de 2.4
GHz y alrededor de 5 GHz. WARDEN trabaja exclusivamente en la banda de 2.4
GHz.

\subsection{¿Por qué WiFi vive en 2.4 GHz?}
La elección no es arbitraria. La banda de 2.4 GHz tiene varias propiedades
útiles:

\begin{itemize}
  \item \textbf{Es banda ISM (Industrial, Scientific and Medical).} Esto
    significa que en la mayoría de países es de uso libre: cualquiera puede
    transmitir en ella sin necesidad de licencia, siempre que respete
    límites de potencia. Otras bandas (por ejemplo, las que usan operadoras
    celulares) están reservadas y se subastan por miles de millones de
    dólares.
  \item \textbf{Penetración razonable de paredes.} Las ondas de 2.4 GHz
    atraviesan paredes domésticas con pérdida moderada. Frecuencias mucho
    más altas (5 GHz, 6 GHz) atraviesan peor; frecuencias mucho más bajas
    requerirían antenas enormes.
  \item \textbf{Antenas pequeñas.} La longitud física de una antena
    eficiente es proporcional a la longitud de onda, que a su vez es
    inversa a la frecuencia. A 2.4 GHz, una antena puede ser de pocos
    centímetros, ideal para dispositivos portátiles.
\end{itemize}

\subsection{La banda de 2.4 GHz no es un solo punto}
Cuando se dice ``WiFi en 2.4 GHz'' en realidad se refiere a una región del
espectro: desde aproximadamente 2.400 GHz hasta 2.495 GHz. Esa región se
divide en \textbf{canales}.

% ==============================================================================
\section{Canales: ¿por qué WiFi solo escucha uno a la vez?}
% ==============================================================================

\subsection{La metáfora de las habitaciones}
Imagina un edificio enorme con muchas habitaciones. En cada habitación hay
gente conversando. Si te paras en el pasillo, oyes un murmullo confuso de
todas las habitaciones a la vez y no entiendes nada. Pero si entras a una
habitación específica, oyes claramente la conversación de esa habitación
y solo esa.

La banda de 2.4 GHz funciona igual. Está dividida en \textbf{14 canales}
numerados del 1 al 14, cada uno centrado en una frecuencia ligeramente
distinta. Cada canal es como una habitación: las redes WiFi que operan en
el canal 6 no interfieren mucho con las que operan en el canal 11, porque
son frecuencias diferentes. Un dispositivo WiFi puede sintonizarse a un
canal específico, y mientras está sintonizado ahí, solo escucha lo que
ocurre en ese canal.

\subsection{Tabla de canales en 2.4 GHz}

\begin{datatable}{L{2cm} L{3cm} L{10.3cm}}{Canal & Frecuencia central & Notas}
\drow{1   & 2412 MHz & Disponible en todo el mundo}
\drow{2   & 2417 MHz & }
\drow{3   & 2422 MHz & }
\drow{4   & 2427 MHz & }
\drow{5   & 2432 MHz & }
\drow{6   & 2437 MHz & Uno de los canales ``no superpuestos'' más usados}
\drow{7   & 2442 MHz & }
\drow{8   & 2447 MHz & }
\drow{9   & 2452 MHz & }
\drow{10  & 2457 MHz & }
\drow{11  & 2462 MHz & Uno de los canales ``no superpuestos'' más usados}
\drow{12  & 2467 MHz & No disponible en Estados Unidos}
\drow{13  & 2472 MHz & No disponible en Estados Unidos}
\drow{14  & 2484 MHz & Solo Japón}
\end{datatable}

\subsection{¿Por qué se recomiendan los canales 1, 6 y 11?}
Aunque los canales son 14, cada uno usa un ancho de banda de 20 MHz, mientras
que la separación entre canales contiguos es solo de 5 MHz. Esto significa
que los canales se \textbf{solapan} entre sí. Solo los canales 1, 6 y 11
están suficientemente separados para no interferirse mutuamente. Por eso son
los más utilizados.

Para WARDEN se utilizará un canal único (a definir cuando se configure el
router de laboratorio), preferentemente alguno con baja densidad de redes
vecinas.

\subsection{Implicación crítica para WARDEN}
Como un dispositivo WiFi solo puede estar sintonizado en un canal a la vez,
nuestro adaptador detector (el Panda) solo escucha lo que ocurre en el canal
donde lo configuremos. Si un ataque sucede en el canal 6 pero el Panda está
en el canal 11, no lo captura. Por eso, antes de cualquier sesión de
laboratorio, el Panda se configura explícitamente en el mismo canal del
router de laboratorio.

% ==============================================================================
\section{Direcciones MAC: la ``cédula'' de cada dispositivo}
% ==============================================================================

\subsection{Cada interfaz de red tiene una identidad}
Toda interfaz de red, ya sea cableada o inalámbrica, tiene asignada una
dirección única llamada \textbf{dirección MAC} (Media Access Control). Es
una secuencia de 6 bytes representada como 6 grupos de dos dígitos
hexadecimales:

\begin{verbatim}
4a:c9:7b:ce:34:46
\end{verbatim}

Esta dirección viene grabada de fábrica en cada chip de red. Es como la
cédula de ciudadanía del dispositivo en la red local: ningún otro chip en
el mundo debería tener exactamente esa misma MAC.

\subsection{El fabricante está codificado en la MAC}
Los primeros 3 bytes de la MAC se llaman \textbf{OUI} (Organizationally
Unique Identifier) y identifican al fabricante del chip. Los últimos 3
bytes son únicos por dispositivo dentro de ese fabricante. Por ejemplo:

\begin{itemize}
  \item Una MAC que comienza con \texttt{00:c0:ca:...} indica un chip
    fabricado por una compañía registrada en IEEE bajo ese OUI.
  \item Existen bases de datos públicas (mantenidas por IEEE) que permiten
    consultar el fabricante a partir del OUI.
\end{itemize}

\subsection{BSSID: la MAC del Access Point}
Cuando hablamos específicamente de la dirección MAC de un \emph{Access
Point}, le damos un nombre más técnico: \textbf{BSSID} (Basic Service Set
Identifier). El BSSID es simplemente la MAC del AP. Cuando un cliente se
conecta a una red WiFi, internamente lo que hace es conectarse a un BSSID
específico.

\subsection{SSID vs BSSID: dos conceptos distintos}
Esta es una distinción que confunde mucho al principio, pero es
fundamental:

\begin{datatable}{L{2cm} L{5cm} L{8.8cm}}{Término & Qué es & Ejemplo}
\drow{SSID  & Nombre legible de la red WiFi (lo que el usuario ve y elige) & \texttt{LAB\_WARDEN\_UTP}}
\drow{BSSID & Dirección MAC del AP (lo que internamente identifica al equipo) & \texttt{e4:ab:89:d6:9b:80}}
\end{datatable}

Una red puede tener \textbf{varios APs con el mismo SSID pero distinto
BSSID} (por ejemplo, una empresa con múltiples routers que se anuncian con
el mismo nombre para que los empleados no tengan que reconectarse al
moverse). Lo contrario también puede pasar: un AP atacante puede crear una
red con el mismo SSID que una legítima pero con BSSID distinto. Esto
último es la base del ataque \textbf{Evil Twin}, que veremos más adelante.

% ==============================================================================
\section{Anatomía de un frame 802.11}
% ==============================================================================

\subsection{Un frame es un mensaje empaquetado}
Cada vez que un dispositivo WiFi transmite algo por aire, lo hace en
forma de un \textbf{frame}. Un frame es un mensaje empaquetado con una
estructura precisa: tiene un encabezado con metadatos (origen, destino,
tipo de mensaje, número de secuencia), un cuerpo con el contenido, y una
verificación final para detectar errores.

Es análogo a un sobre postal: el sobre tiene remitente, destinatario,
estampilla, sello, y dentro va la carta. Si llega dañado, se nota.

\subsection{Tres categorías de frames}
El estándar 802.11 define tres categorías principales de frames, cada una
con un propósito distinto:

\begin{datatable}{L{3.5cm} L{12.3cm}}{Categoría & Propósito y ejemplos}
\drow{Frames de gestión & Coordinan el uso de la red. No transportan datos del usuario, sino información sobre la red en sí: anuncios de APs, solicitudes de conexión, desconexiones, etc. \emph{Ejemplos: Beacon, Probe Request, Probe Response, Authentication, Association Request, Deauthentication.}}
\drow{Frames de control & Sincronizan transmisiones y confirman recepciones. Son muy cortos y muy frecuentes. \emph{Ejemplos: ACK (acknowledgment), RTS (request to send), CTS (clear to send).}}
\drow{Frames de datos   & Transportan información real del usuario: páginas web, mensajes, video, audio, etc. Aquí va el contenido cifrado por WPA2.}
\end{datatable}

\textbf{WARDEN se centra casi exclusivamente en frames de gestión}, porque
ahí es donde se ven los ataques que estudiamos. Los beacons, los deauth,
los probe responses son todos frames de gestión.

\subsection{Por qué un cliente normal no ve la mayoría de los frames}
Un dispositivo en operación normal opera en modo \emph{managed} (lo
veremos a fondo más adelante). En este modo, la tarjeta WiFi descarta
casi todos los frames que recibe del aire, salvo aquellos dirigidos
específicamente a su MAC o a una dirección broadcast relevante. La
mayoría de frames que pasan por el aire (beacons de redes vecinas, ACKs
ajenos, datos de otros dispositivos) son ignorados antes de que el sistema
operativo los vea.

Esto es eficiente para uso normal, pero hace imposible la detección de
ataques. Un atacante puede emitir cientos de frames maliciosos y un
cliente normal ni se entera. Para verlos, hay que cambiar el modo de
operación de la tarjeta, lo que nos lleva al concepto siguiente.

% ==============================================================================
\section{Modos de operación de una tarjeta WiFi}
% ==============================================================================

\subsection{La misma tarjeta, distintos modos}
Una tarjeta WiFi puede operar en varios modos, configurables por software.
Los más relevantes son:

\begin{datatable}{L{3cm} L{12.8cm}}{Modo & Comportamiento}
\drow{Managed   & Modo cliente. La tarjeta se asocia a un AP y solo recibe frames dirigidos a ella. Es el modo por defecto en cualquier laptop o celular.}
\drow{AP        & La tarjeta actúa como Access Point: emite beacons, acepta conexiones de clientes. Es el modo de un router WiFi o un hotspot de celular.}
\drow{Monitor   & Modo ``oyente pasivo''. La tarjeta no se asocia a ninguna red y captura \emph{todos} los frames que recibe en el canal sintonizado, sin filtrar por destinatario.}
\drow{Mesh, IBSS, P2P & Modos para arquitecturas específicas (redes ad-hoc, redes mesh, conexiones directas tipo Wi-Fi Direct). No los usamos en WARDEN.}
\end{datatable}

\subsection{Modo monitor: el modo ``Dios'' para análisis}
El modo monitor convierte a la tarjeta en un sensor pasivo del aire. En este
modo, la tarjeta:

\begin{itemize}
  \item No se conecta a ninguna red.
  \item No transmite nada (en operación pura, aunque algunos drivers
    permiten también inyección).
  \item Recibe \emph{cada frame} 802.11 que llega físicamente al
    adaptador, sin importar el destinatario.
  \item Entrega esos frames al sistema operativo con todos sus metadatos
    intactos: encabezado de capa de enlace, información de radiofrecuencia
    (potencia de señal recibida, canal exacto, modulación usada), etc.
\end{itemize}

Es la capacidad fundamental que hace posible el subsistema defensivo de
WARDEN. Sin modo monitor, no podemos ver beacons falsos, no podemos ver
deauths dirigidos a otros, no podemos detectar nada.

\subsection{No todos los adaptadores soportan modo monitor}
Aquí hay un detalle crítico: \textbf{la mayoría de adaptadores WiFi en
laptops y celulares no soportan modo monitor en Linux}, no porque el chip
no pueda hacerlo físicamente, sino porque el driver del fabricante no lo
expone. Los fabricantes desactivan deliberadamente esta capacidad por
varias razones (simplificación, supuestas mejoras de seguridad, requisitos
regulatorios mal interpretados).

Por eso, en proyectos como WARDEN se utilizan adaptadores específicos cuyos
drivers libres soportan bien el modo monitor. El \textbf{Panda Wireless
PAU0B AC600}, con chipset MediaTek MT7610U, es uno de esos adaptadores. Su
driver Linux (\texttt{mt76x0u}) es de código abierto y permite modo monitor
estable en la banda 2.4 GHz.

\subsection{Implicación para el laboratorio WARDEN}
La estación detectora del proyecto necesita el adaptador Panda
\emph{específicamente porque} la WiFi interna de la laptop, en la mayoría de
casos, no soporta modo monitor. El Panda no es un capricho ni una mejora de
rendimiento: es una pieza imprescindible de hardware sin la cual el
proyecto no es posible.

% ==============================================================================
\section{Cómo se conecta un cliente a una red: la asociación}
% ==============================================================================

\subsection{Pasos de una conexión WiFi normal}
Cuando una laptop o celular ``se conecta a una WiFi'', el proceso interno
involucra varios pasos que el usuario no ve:

\begin{enumerate}
  \item \textbf{Descubrimiento.} El cliente escucha beacons o emite
    \emph{probe requests} preguntando ``¿está la red X cerca?''. Los APs
    responden con \emph{probe responses}.
  \item \textbf{Autenticación.} El cliente envía un frame de autenticación.
    En redes modernas este paso es trivial (siempre se aprueba); el
    ``verdadero'' control de acceso ocurre después con WPA2.
  \item \textbf{Asociación.} El cliente envía una \emph{Association
    Request} al AP. Si el AP acepta, responde con una \emph{Association
    Response} y desde ese momento el cliente está ``asociado'' a esa red.
  \item \textbf{Handshake WPA2 (de 4 vías).} En redes con WPA2, después de
    la asociación se ejecuta un intercambio de 4 mensajes que establece la
    clave de cifrado de la sesión. Sin este paso, el cliente está asociado
    pero no puede transmitir datos.
  \item \textbf{Tráfico cifrado.} A partir de aquí, todos los frames de
    datos van cifrados con la clave establecida en el handshake.
\end{enumerate}

\subsection{Por qué importa entender los pasos}
Cada uno de estos pasos involucra frames específicos visibles en el aire.
El detector de WARDEN observa estos frames para distinguir entre tráfico
normal y patrones anómalos asociados a ataques. Por ejemplo, un cliente
normal emite \emph{probe requests} de manera intermitente; un cliente bajo
ataque de deauth los emite repetidamente porque está intentando reconectarse.

% ==============================================================================
\section{El handshake WPA2 y por qué un atacante quiere capturarlo}
% ==============================================================================

\subsection{¿Qué es el handshake?}
El \textbf{handshake de 4 vías} (en inglés \emph{4-way handshake}) es un
intercambio de cuatro mensajes que ocurre cada vez que un cliente se conecta
a una red WPA2. El propósito es que ambos extremos (cliente y AP) deriven
una clave de cifrado de sesión sin nunca enviar la contraseña real por el
aire.

El nombre técnico de este protocolo es \textbf{4-Way Handshake EAPOL}. Los
mensajes contienen valores aleatorios, identificadores, y un código MIC
(Message Integrity Code) que prueba que ambos extremos conocen la
contraseña, sin exponerla.

\subsection{¿Por qué un atacante captura el handshake?}
Aunque la contraseña real nunca viaja por el aire, los mensajes del
handshake contienen información derivada de ella. Un atacante que captura
los cuatro mensajes puede luego, \textbf{offline} (sin conectividad ni
interacción con la red), probar contraseñas candidatas a alta velocidad
hasta encontrar una que produzca el mismo MIC. Esto es el llamado ``ataque
de diccionario offline''.

\textbf{El handshake es el activo principal en ataques contra WPA2.} Si un
atacante lo captura y la contraseña es débil (común, predecible, en algún
diccionario público), puede recuperar la contraseña en minutos u horas.

\subsection{¿Cómo se captura el handshake?}
El handshake solo ocurre en el momento de la conexión. Hay dos formas de
capturarlo:

\begin{itemize}
  \item \textbf{Esperar pasivamente} a que un cliente se conecte
    naturalmente (pueden ser horas o días).
  \item \textbf{Forzar la captura} desconectando a un cliente ya conectado
    para que se reconecte de inmediato. Esto se hace con un ataque de
    \textbf{deauthentication}, que veremos a continuación.
\end{itemize}

\subsection{Aclaración importante para WARDEN}
El proyecto WARDEN \emph{no incluye} cracking offline de contraseñas. La
captura del handshake aparece en el documento de laboratorio solo como
ejercicio de práctica para entender cómo se ven los deauths en el aire.
WARDEN se enfoca en \emph{detectar} los ataques, no en explotarlos hasta el
final.

% ==============================================================================
\section{Beacons: el ``latido'' de cada red WiFi}
% ==============================================================================

\subsection{¿Qué es un beacon?}
Un \textbf{beacon} (en español ``baliza'') es un frame de gestión que cada
Access Point emite \textbf{periódicamente} para anunciar su existencia. La
frecuencia típica es cada 102.4 milisegundos, lo que da aproximadamente 10
beacons por segundo por cada AP activo.

Cada beacon contiene:

\begin{itemize}
  \item \textbf{SSID:} el nombre de la red.
  \item \textbf{BSSID:} la MAC del AP.
  \item \textbf{Capacidades:} qué cifrados soporta (WPA2, WPA3, abierto,
    etc.), qué velocidades, etc.
  \item \textbf{Información de canal y banda.}
  \item \textbf{Marca de tiempo y número de secuencia.}
\end{itemize}

\subsection{¿Para qué sirven los beacons?}
Los beacons son lo que permite que tu celular o laptop muestre una lista
de redes WiFi disponibles. Cuando abres la configuración WiFi y ves ``Mi
Red'', ``Vecino123'', ``Cafetería UTP'', lo que estás viendo es la información
extraída de los beacons que esas redes están emitiendo continuamente.

\subsection{¿Qué es el Beacon Flooding?}
El \textbf{Beacon Flooding} (inundación de beacons) es un ataque que
consiste en \textbf{emitir miles de beacons falsos por segundo}, cada uno
anunciando una red distinta con SSID y BSSID inventados. El efecto es:

\begin{itemize}
  \item La lista de redes WiFi de los dispositivos cercanos se llena de
    redes falsas.
  \item Las herramientas automatizadas que escanean redes (incluidos
    sistemas de seguridad) tienen dificultades para distinguir señal real
    de ruido.
  \item En casos extremos, algunos clientes presentan errores o lentitud.
\end{itemize}

Es la primera fase de la cadena ofensiva de WARDEN. Su firma de detección
es relativamente clara: una tasa de beacons únicos por segundo muy
superior a lo normal, todos provenientes del mismo emisor físico (la
ESP32).

% ==============================================================================
\section{Deauthentication: el ``cuelgue forzado''}
% ==============================================================================

\subsection{¿Qué es un frame de deauthentication?}
Un frame de \textbf{deauthentication} (abreviado \textbf{deauth}) es un
frame de gestión que indica el fin de una sesión de autenticación. Cuando
un cliente se desconecta voluntariamente de una red, su sistema operativo
envía un deauth al AP. Cuando el AP necesita expulsar a un cliente (por
ejemplo, por sobrecarga), envía un deauth al cliente.

En operación normal, los deauths son raros: solo aparecen cuando alguien
cierra la WiFi, sale de cobertura, o se reinicia un router.

\subsection{El gran problema de los deauths}
Aquí está la vulnerabilidad histórica de 802.11: \textbf{los frames de
deauth no estaban autenticados}. Esto significa que cualquiera puede
fabricar un deauth haciéndose pasar por el AP legítimo, y el cliente lo
acepta sin cuestionar. Una vez recibido, el cliente se desconecta
inmediatamente.

Un atacante con un adaptador WiFi y conocimiento básico puede emitir
deauths a voluntad y desconectar a cualquier cliente cercano de su red.

\subsection{¿Para qué sirve un ataque de deauth?}
El deauth raramente es un fin en sí mismo. Su utilidad es habilitar otros
ataques:

\begin{itemize}
  \item \textbf{Forzar reconexión}, para capturar el handshake WPA2.
  \item \textbf{Forzar reconexión}, para que el cliente se conecte a un
    Evil Twin que el atacante tiene preparado (esta es la fase 2 de la
    cadena WARDEN).
  \item \textbf{Denegación de servicio}, manteniendo a una víctima
    desconectada por tiempo prolongado.
\end{itemize}

\subsection{La mitigación: 802.11w (PMF)}
En 2009 se publicó la enmienda 802.11w, también llamada \textbf{Protected
Management Frames (PMF)}, que añade autenticación a los frames de gestión
(incluidos los deauths). Cuando un cliente y un AP soportan PMF, los
deauths falsificados son ignorados.

PMF es obligatorio en WPA3 pero opcional en WPA2. La mayoría de equipos
domésticos modernos lo soportan pero no lo activan por defecto. Para
WARDEN, el router de laboratorio se configurará sin PMF para que los
ataques sean reproducibles.

% ==============================================================================
\section{Evil Twin: el ``gemelo malvado''}
% ==============================================================================

\subsection{La idea del ataque}
Un \textbf{Evil Twin} es un Access Point malicioso que imita el SSID de
una red legítima con la intención de engañar a clientes para que se
conecten a él en lugar de al AP real. Una vez conectado al Evil Twin, el
atacante controla todo el tráfico del cliente: puede observarlo,
modificarlo, redirigirlo, o simplemente capturar las credenciales que el
cliente intente enviar.

\subsection{¿Por qué funciona?}
La razón principal es que los dispositivos WiFi modernos (celulares,
laptops) están configurados para \textbf{reconectarse automáticamente} a
redes conocidas. Cuando un dispositivo ve un SSID al que se ha conectado
antes, intenta conectarse sin pedir confirmación al usuario.

Si el Evil Twin se anuncia con el mismo SSID que una red conocida y, por
algún motivo, el cliente prefiere conectarse a él (por ejemplo, porque la
red real fue desconectada con un deauth), el cliente cae en la trampa sin
darse cuenta.

\subsection{Variantes del Evil Twin}
Existen dos variantes principales:

\begin{itemize}
  \item \textbf{Evil Twin abierto:} El AP malicioso se anuncia como red
    abierta (sin contraseña). Funciona porque muchos clientes, al ver una
    red conocida abierta, se conectan asumiendo que es una versión ``guest''
    legítima.
  \item \textbf{Evil Twin protegido:} El AP malicioso replica también la
    configuración de seguridad de la red real. Es más difícil de
    implementar pero más convincente.
\end{itemize}

WARDEN implementa la variante \textbf{abierta}, que es la más usada en
ataques reales y la más fácil de demostrar.

\subsection{El portal cautivo: el cierre del ataque}
Una vez que la víctima se conecta al Evil Twin, el atacante necesita
capturar algo de valor. La técnica clásica es presentar un \textbf{portal
cautivo falso}: una página web que aparece automáticamente al conectarse
a una red abierta, normalmente para ``iniciar sesión'' o ``aceptar términos''.

Los usuarios están acostumbrados a ver portales cautivos en aeropuertos,
hoteles y cafeterías, así que un portal falso bien hecho no levanta
sospechas. Cuando el usuario ingresa sus credenciales (correo,
contraseña, número de teléfono), estas son capturadas por el atacante.

En WARDEN, el portal cautivo es una página web minimalista servida por la
propia ESP32, que registra los datos ingresados y los muestra por consola
serial como evidencia de captura exitosa. Las credenciales utilizadas
durante las pruebas son siempre ficticias.

% ==============================================================================
\section{La cadena de ataque WARDEN: las tres fases en orden}
% ==============================================================================

\subsection{Por qué las tres fases tienen sentido juntas}
Hasta aquí hemos visto los tres ataques por separado. Lo crítico para
entender WARDEN es que en la práctica real \emph{rara vez se ejecutan
solos}: se encadenan en una secuencia que multiplica su efectividad. La
cadena completa es:

\begin{enumerate}
  \item \textbf{Beacon Flooding (Reconocimiento y Distracción).} El
    atacante satura el entorno con beacons falsos. La víctima ve
    decenas de redes nuevas en su lista; el ruido reduce su capacidad de
    distinguir lo legítimo de lo sospechoso.
  \item \textbf{Deauthentication (Aislamiento).} El atacante desconecta a
    la víctima de su red real con frames de deauth falsificados. La
    víctima pierde su conexión y comienza a buscar redes para
    reconectarse.
  \item \textbf{Evil Twin (Explotación).} El atacante presenta un AP
    abierto con el SSID de la red real. La víctima, configurada para
    reconectarse automáticamente y viendo su red conocida disponible, se
    conecta al Evil Twin. El portal cautivo falso captura sus
    credenciales.
\end{enumerate}

\subsection{Por qué esta cadena es la firma del proyecto}
WARDEN no detecta tres ataques aislados: detecta una \emph{cadena} con un
patrón temporal específico. Tres firmas en intervalo corto (Beacon Flood,
luego Deauth, luego Evil Twin) son una evidencia mucho más fuerte de un
ataque coordinado que cualquiera de las firmas por separado. Esto es lo
que hace que el proyecto sea conceptualmente más interesante que un IDS
simple basado en una sola regla.

% ==============================================================================
\section{Cierre y mapa de lectura siguiente}
% ==============================================================================

Si llegaste hasta aquí entendiendo la mayor parte, ya tienes el modelo
mental que necesitas para participar en WARDEN. Los términos y conceptos
clave que ahora deberías conocer:

\begin{itemize}
  \item Qué es WiFi, qué es 802.11, y por qué se trabaja en 2.4 GHz.
  \item Qué son los canales y por qué un adaptador solo puede oír uno a
    la vez.
  \item Qué es una dirección MAC, qué es BSSID, qué es SSID y la
    diferencia entre ambos.
  \item Qué es un frame y los tres tipos principales (gestión, control,
    datos).
  \item Qué es modo monitor y por qué necesitamos el adaptador Panda.
  \item Qué es un beacon, un deauth, una asociación, un handshake.
  \item Qué es Beacon Flooding, qué es un ataque de Deauth, qué es un
    Evil Twin, y por qué los tres encadenados son más potentes que
    cualquiera por separado.
\end{itemize}

\subsection{Siguiente lectura}
Con estos fundamentos asimilados, el siguiente documento que conviene leer
es \emph{WARDEN - Manual de Laboratorio}, que toma todos los conceptos
explicados aquí y los aterriza en comandos concretos sobre Linux, paso a
paso, para que puedas reproducir los tres ataques de forma manual antes
de que los automaticemos en la ESP32.

El manual de laboratorio asume que ya entendiste este documento. Si en
algún momento durante la práctica algo no tiene sentido, lo más probable
es que la respuesta esté aquí.

\end{document}
